\documentclass[a5paper,10pt]{article}

% https://github.com/InDevRus/filippov-solutions

% Нумерация страниц
\pagenumbering{gobble}

% Настройка полей
\usepackage{geometry}
\geometry{left=1cm}
\geometry{right=1cm}
\geometry{top=1cm}
\geometry{bottom=1.5cm}

% Вставка таблицы
\usepackage{array}
\usepackage{tabularx}

% Инструменты выравнивания формул
\usepackage{amsmath}
\usepackage{mathtools}

% Диакритика в математике
\usepackage{amsfonts}

% Вычисления размеров на ходу
\usepackage{calc}

% Удобные записи производных
\usepackage[italic=true]{derivative}

% Настройки сносок
\usepackage{enotez}

% Длинные знаки векторов
\usepackage{anyfontsize}
\usepackage[b]{esvect}

% Вставка картинок
\usepackage{graphicx}

% Вставка красной строки
\usepackage{indentfirst}
\setlength{\parindent}{1em}

% Условия в командах
\usepackage{xifthen}

% Луа-скрипты
\usepackage{luacode}

% Настройка команд отдельно в математическом и текстовом режимах
\usepackage{mathcommand}

% Для повторения LaTeX кода
\usepackage{multido}

% Конфигурация отступов формул
\delimitershortfall-1sp
\usepackage{mleftright}
\mleftright

% Растягивание объектов
\usepackage{relsize}

% Обтекание текстом
\usepackage{wrapfig}
\usepackage{subcaption}
\usepackage{floatflt}

% Поддержка ввода кириллицы
\usepackage[english, russian]{babel}

% Настройка корневой позиции для компилятора
\usepackage{import}
\providecommand*{\ProjectRootPath}{..}

\import{\ProjectRootPath/preambles/macros/}{theorems}

\import{\ProjectRootPath/preambles/fonts}{font_setup}

% Знак градуса
\usepackage{siunitx}

\import{\ProjectRootPath/preambles/macros/}{operators}
\import{\ProjectRootPath/preambles/macros/}{redefinitions}

\import{\ProjectRootPath/preambles/fonts}{letters_setup}
\import{\ProjectRootPath/preambles/fonts/adjustments}{custom_superscripts}

\import{\ProjectRootPath/preambles/custom}{formatting}
\import{\ProjectRootPath/preambles/custom}{frames}
\import{\ProjectRootPath/preambles/custom}{list_setup}

% TODO: перевести команды на современный стиль объявления

\begin{document}

В обозначениях задачи \textbf{737}
$\psi\br{x} = \pow{x}{\selevate{-\frac {1}{2}}}$,
$\phi\`\br{x} = \pow{\psi}{-2}\br{x} = x$,
например,
$\phi\br{x} = \frac{1}{2} \pow{x}{2}$.
Тогда заменяем \vspace{0.4em}
\begin{align*}
    \tsys{
        t = \phi\br{x} = \frac{1}{2} \pow{x}{2} \\
        y = u \psi\br{x} = u \pow{x}{-\frac {1}{2}}
    }
    &&
    \sub{Q}{u\br{t}}\br{t} 
    = -1 + \pow{\psi}{3}\br{x} \psi\`\`\br{x}
    = -1 + \dfrac{3}{4} \pow{x}{-4}
    = -1 + \dfrac{3}{16} \pow{t}{-2},
\end{align*}
в результате чего получаем 
$u\`\` + \br{-1 + \dfrac {3}{16}\pow{t}{-2}}\,u\br{t} = 0$.
При $t \to +\infty$ результат таков:
\begin{align*}
    \sub{u}{1}\br{t} = \pow{e}{t} \br{1 + O\br{\pow{t}{-1}}}
    \qquad
    \sub{u}{2}\br{t} = \pow{e}{-t} \br{1 + O\br{\pow{t}{-1}}}
\end{align*}
\begin{align*}
    \sub{y}{1}\br{x} 
    = \pow{x}{\selevate{-\frac {1}{2}}}
    \pow{e}{\selevate{\frac {1}{2} \pow{x}{2}}}
    \br{1 + O\br{\pow{x}{-2}}}
    \qquad
    \sub{y}{2}\br{x} 
    = \pow{x}{\selevate{-\frac {1}{2}}}
    \pow{e}{\selevate{-\frac {1}{2} \pow{x}{2}}}
    \br{1 + O\br{\pow{x}{-2}}}
\end{align*}

\end{document}
